% Chapter 2 — Background

\chapter{Background}
\label{Chapter2}

%----------------------------------------------------------------------------------------

\section{Missing Data Mechanisms}

The behaviour of any imputation method depends critically on the mechanism that produces the missing values. Rubin (1976) established the canonical taxonomy of three missing data mechanisms, which determines the conditions under which different imputation strategies are valid \cite{Rubin1976}.

\textbf{Missing Completely At Random (MCAR).} A value is MCAR if the probability of its being missing is independent of both the observed data and the missing data. Formally, if $R$ is the missingness indicator matrix ($R_{ij} = 1$ if $X_{ij}$ is observed, 0 if missing), then MCAR requires $P(R \mid X) = P(R)$. Under MCAR, complete-case analysis is unbiased, though it discards information. MCAR is the most restrictive assumption and rarely holds in practice.

\textbf{Missing At Random (MAR).} A value is MAR if the probability of its being missing may depend on the observed data but not on the missing data itself: $P(R \mid X) = P(R \mid X_{\text{obs}})$. Under MAR, the mechanism is ignorable \cite{Rubin1976} --- valid imputation can be performed using only the observed data, without modelling the missingness mechanism. Most modern imputation methods, including MICE and MIGA, assume MAR.

\textbf{Missing Not At Random (MNAR).} A value is MNAR if its probability of being missing depends on the value itself, even after conditioning on observed data: $P(R \mid X) \neq P(R \mid X_{\text{obs}})$. MNAR is the most challenging case. Examples include income values that are more likely to be missing for high earners, or medical test results missing for patients who did not attend follow-up appointments. Under MNAR, the missingness mechanism is non-ignorable: any imputation method that ignores the mechanism will produce biased results. Chapter~\ref{Chapter4} provides the first empirical characterisation of MIGA's failure mode under MNAR.

The three mechanisms form a hierarchy of restrictiveness. MCAR implies MAR, and MAR implies the mechanism is ignorable. Van Buuren (2018, \S1.2) provides an extensive treatment of how to reason about mechanisms in practice \cite{vanBuuren2018}.

%----------------------------------------------------------------------------------------

\section{Classical Imputation Methods}

\subsection{Complete-Case Analysis}

Complete-case analysis (listwise deletion) discards any observation with at least one missing value, performing the analysis on the complete cases only. Under MCAR, this produces unbiased estimates but reduces statistical power. Under MAR or MNAR, it introduces selection bias: the complete cases are a non-representative subset of the population. With 30\% missing values spread across features, the fraction of complete rows can fall below 10\% even for moderate $p$. Complete-case analysis is therefore rarely appropriate for real-world datasets.

\subsection{Single Imputation}

Single imputation replaces each missing value with a single plausible estimate. The simplest method substitutes the column mean (for continuous features) or the column mode (for categorical features). Mean imputation is computationally trivial and introduces no bias in the mean itself, but it severely distorts the variance structure: replacing 30\% of a column with its mean compresses the variance by approximately $(1 - 0.3)^2 \approx 0.49$, nearly halving it.

KNN imputation \cite{Troyanskaya2001} replaces each missing value with a weighted average of the $k$ nearest complete observations in the feature space. KNN preserves local structure better than mean imputation and is non-parametric, but its quality depends on the choice of $k$ and the distance metric, and it does not account for the uncertainty introduced by imputation.

Both mean and KNN imputation are single-imputation methods: they produce one complete dataset and treat imputed values as if they were observed. This understates the uncertainty due to missingness and produces overconfident inferences \cite{Rubin1987}.

\subsection{Multiple Imputation by Chained Equations (MICE)}

Multiple imputation generates $M > 1$ complete datasets by drawing imputed values from the predictive distribution $P(X_{\text{mis}} \mid X_{\text{obs}})$, combines the analysis results across datasets using Rubin's rules, and produces inferences that properly account for missingness uncertainty \cite{Rubin1987}.

MICE \cite{vanBuuren2011}, also known as Fully Conditional Specification (FCS), is the dominant multiple imputation method. For a dataset with $p$ partially observed variables, MICE iterates a cycle of $p$ conditional imputation models: for each variable $j$, it fits a regression model of $X_j$ on all other variables $X_{-j}$ using the currently completed data, then draws imputations from the posterior predictive distribution of this model. The cycle repeats until convergence.

The key strength of MICE is its flexibility: each variable can use a different model family (linear regression for continuous variables, logistic regression for binary variables). MICE is theoretically valid under MAR and produces imputations with the correct conditional distribution of each variable given the others.

However, MICE's optimality criterion is minimum expected squared error, which corresponds to imputing conditional means. Van Buuren (2018, \S2.6) explicitly states:

\begin{quotation}
``The minimum mean squared error is achieved by regression imputation, which suppresses variance and produces biased statistical inference.''
\end{quotation}

This is a structural property: by imputing $\mathbb{E}[X_j \mid X_{-j}]$, MICE systematically underestimates the marginal variance $\text{Var}(X_j)$. By the law of total variance:
\begin{equation}
    \text{Var}(X_j) = \mathbb{E}[\text{Var}(X_j \mid X_{-j})] + \text{Var}(\mathbb{E}[X_j \mid X_{-j}])
\end{equation}
MICE discards the first term. For analyses that require the correct marginal distribution, MICE produces biased results. This is the fundamental motivation for MIGA.

%----------------------------------------------------------------------------------------

\section{The \texorpdfstring{$F_r$}{Fr}--RMSE Orthogonality}
\label{sec:orthogonality}

A central claim of this thesis is that distributional fidelity (measured by $F_r$) and pointwise accuracy (measured by RMSE) are \emph{orthogonal} objectives: optimising one systematically degrades the other. This section formalises that claim, first by establishing a general mathematical framework, then by proving the orthogonality as a consequence.

\subsection{Mathematical Framework}
\label{sec:orthogonality_framework}

\textbf{Setup.} Let $X \in \mathbb{R}^{n \times p}$ be a dataset with $n$ observations and $p$ features. Let $M \in \{0,1\}^{n \times p}$ be the missingness mask ($M_{ij} = 0$ if $X_{ij}$ is missing). Under MAR, $M \perp X_{\mathrm{mis}} \mid X_{\mathrm{obs}}$. An \emph{imputation map} is any measurable function $\mathcal{F}: (X_{\mathrm{obs}},\, M) \mapsto \hat{X}_{\mathrm{mis}}$.

\begin{definition}[Irreducible conditional variance]
\label{def:tau}
For feature $j$ with at least one missing value, define the \emph{irreducible conditional variance}:
\begin{equation}
    \tau_j \;=\; \mathbb{E}\!\bigl[\mathrm{Var}(X_j \mid X_{-j})\bigr]
              \;=\; \mathrm{Var}(X_j)\,(1 - R_j^2),
    \label{eq:tau_j}
\end{equation}
where $R_j^2$ is the population coefficient of determination of $X_j$ regressed on $X_{-j}$.
Write $\tau = \sum_{j \in \mathcal{J}} \tau_j$ where $\mathcal{J}$ is the set of features with missing values.
\end{definition}

$\tau_j$ is the variance of $X_j$ that cannot be recovered from the remaining features: it is the \emph{aleatoric} uncertainty of imputation.  When $R_j^2 \approx 1$ (feature $j$ is nearly linearly predictable from $X_{-j}$), $\tau_j \approx 0$ and RMSE-optimal imputation is nearly distributional-faithful; when $R_j^2 \approx 0$ (feature $j$ is independent of $X_{-j}$), $\tau_j = \mathrm{Var}(X_j)$ and the deficiency is maximal.

\textbf{RMSE decomposition.} For any imputation map $\mathcal{F}$, the bias--variance decomposition of squared error gives:
\begin{align}
    \mathbb{E}\!\bigl[\mathrm{RMSE}(\mathcal{F})^2\bigr]
    &= \underbrace{\mathbb{E}\!\bigl[\|\mathcal{F} - \mathbb{E}[X_{\mathrm{mis}} \mid X_{\mathrm{obs}}]\|^2\bigr]}_{\text{estimation error} \;\geq\; 0}
    + \underbrace{\tau}_{\text{irreducible variance}}.
    \label{eq:rmse_decomp}
\end{align}
The estimation error is non-negative and vanishes only when $\mathcal{F} = \mathbb{E}[X_{\mathrm{mis}} \mid X_{\mathrm{obs}}]$.
Therefore the unique RMSE-minimising imputer is:
\begin{equation}
    \mathcal{F}^* \;=\; \mathbb{E}[X_{\mathrm{mis}} \mid X_{\mathrm{obs}}],
    \qquad
    \mathrm{RMSE}^* \;=\; \sqrt{\tau},
    \label{eq:mmse_imputer}
\end{equation}
and no imputer can achieve $\mathrm{RMSE} < \sqrt{\tau}$.

\begin{theorem}[Quantitative $F_r$ lower bound for MMSE imputation]
\label{thm:fr_lower_bound}
Let $\pi_j \in (0,1]$ be the proportion of rows in $X_C$ for which feature $j$ is missing, and let $R_j^2$ be as in Definition~\ref{def:tau}. Under the MMSE imputer $\mathcal{F}^*$:
\begin{equation}
    F_r(\mathcal{F}^*) \;\geq\; D_r\!\bigl(\tilde{S}(\mathcal{F}^*),\, I\bigr)
    \;\geq\; \max_{j \in \mathcal{J}} \;\pi_j\,(1 - R_j^2).
    \label{eq:fr_lower_bound}
\end{equation}
\end{theorem}

\begin{proof}
For $j \in \mathcal{J}$, the completed block $X_C$ under $\mathcal{F}^*$ mixes two sub-populations: the $(1-\pi_j)$ fraction of rows in which $X_j$ is observed (contributing marginal variance $\mathrm{Var}(X_j)$) and the $\pi_j$ fraction in which $X_j$ is imputed via $\mathbb{E}[X_j \mid X_{-j}]$ (contributing variance $R_j^2\,\mathrm{Var}(X_j)$). Under MAR the two sub-populations share the same mean, so the sample variance of column $j$ after MMSE imputation satisfies:
\[
    [S_C(\mathcal{F}^*)]_{jj} \;\approx\; [S_A]_{jj}\,\bigl(1 - \pi_j(1 - R_j^2)\bigr).
\]
The relative covariance $\tilde{S} = S_A^{-1/2} S_C S_A^{-1/2}$ therefore has diagonal entry $\tilde{S}_{jj} \approx 1 - \pi_j(1-R_j^2)$, so $|\tilde{S}_{jj} - 1| = \pi_j(1-R_j^2)$. Since $D_r(\tilde{S},I) \geq |\tilde{S}_{jj} - 1|$ for any $j$, and $F_r \geq D_r(\tilde{S}, I)$, the bound follows.
\end{proof}

\begin{remark}
Theorem~\ref{thm:fr_lower_bound} has a direct empirical interpretation. The bound $\pi_j(1-R_j^2)$ is largest when a feature has both high missingness and low predictability from other features. On Haberman, the Survival feature has $R_j^2 \approx 0.07$ and $\pi_j = 0.30$, giving a bound of $\approx 0.28$; this explains the large empirical $F_r$ gap between MIGA and MICE in Chapter~\ref{Chapter4}. On Glass, strong inter-feature correlations ($R_j^2 \to 1$ for most features) suppress the bound, explaining why MICE achieves competitive $F_r$ despite its MMSE objective --- the variance deficiency is small because the features are nearly deterministically predictable from one another.
\end{remark}

\begin{proposition}[RMSE cost of proper imputation]
\label{prop:proper_rmse}
Let $\hat{X}_{\mathrm{mis}} \sim P(X_{\mathrm{mis}} \mid X_{\mathrm{obs}})$ be a \emph{proper} imputation, drawn independently of the true $X_{\mathrm{mis}}$ (van Buuren, 2018, \S3.4). Then:
\begin{equation}
    \mathbb{E}\!\bigl[\mathrm{RMSE}^2\bigr] \;=\; 2\tau.
    \label{eq:proper_rmse}
\end{equation}
That is, proper imputation achieves exactly $\sqrt{2}$ times the minimum RMSE.
\end{proposition}

\begin{proof}
Given $X_{\mathrm{obs}}$, both $X_{\mathrm{mis}}$ and $\hat{X}_{\mathrm{mis}}$ are independent draws from $P(X_{\mathrm{mis}} \mid X_{\mathrm{obs}})$. Therefore:
\[
    \mathbb{E}\!\bigl[\|X_{\mathrm{mis}} - \hat{X}_{\mathrm{mis}}\|^2 \mid X_{\mathrm{obs}}\bigr]
    = \mathrm{Var}(X_{\mathrm{mis}} \mid X_{\mathrm{obs}}) + \mathrm{Var}(\hat{X}_{\mathrm{mis}} \mid X_{\mathrm{obs}})
    = 2\,\mathrm{Var}(X_{\mathrm{mis}} \mid X_{\mathrm{obs}}),
\]
where the cross-term vanishes by $X_{\mathrm{mis}} \perp \hat{X}_{\mathrm{mis}} \mid X_{\mathrm{obs}}$.
Taking the outer expectation gives $\mathbb{E}[\mathrm{RMSE}^2] = 2\tau$.
\end{proof}

\begin{theorem}[$F_r$--RMSE trade-off bound]
\label{thm:tradeoff_bound}
Let $\tau > 0$. Then:
\begin{enumerate}
    \item[(i)] The minimum achievable RMSE is $\sqrt{\tau}$, attained uniquely by $\mathcal{F}^*$.
    \item[(ii)] At $\mathcal{F}^*$: $F_r \geq \max_{j \in \mathcal{J}}\,\pi_j(1 - R_j^2) > 0$.
    \item[(iii)] Any $\mathcal{F}$ achieving $F_r = 0$ must have $\mathrm{RMSE} > \sqrt{\tau}$.
    \item[(iv)] The point $(F_r,\,\mathrm{RMSE}) = (0,\,\sqrt{\tau})$ is unachievable.
\end{enumerate}
\end{theorem}

\begin{proof}
Parts (i) and (ii) follow from Eq.~\eqref{eq:mmse_imputer} and Theorem~\ref{thm:fr_lower_bound}. For (iii): from part (ii), $F_r(\mathcal{F}^*) > 0$, so any $\mathcal{F}$ with $F_r(\mathcal{F}) = 0$ must differ from $\mathcal{F}^*$. By the RMSE decomposition \eqref{eq:rmse_decomp}, any $\mathcal{F} \neq \mathcal{F}^*$ has $\mathrm{RMSE}(\mathcal{F})^2 = \|\mathcal{F} - \mathcal{F}^*\|^2 + \tau > \tau$, so $\mathrm{RMSE}(\mathcal{F}) > \sqrt{\tau}$. Part (iv) follows from (i) and (iii). \end{proof}

\begin{remark}
Proposition~\ref{prop:proper_rmse} shows that proper imputation --- the ideal distributional-fidelity imputer --- pays a factor of $\sqrt{2}$ in RMSE relative to the minimum. MIGA approximates proper imputation via the GA: its RMSE penalty over MICE (1.4--2.1$\times$ empirically; Chapter~\ref{Chapter4}) is consistent with this $\sqrt{2}$ theoretical bound, with the excess arising from imperfect distributional matching.
\end{remark}

Theorems~\ref{thm:fr_lower_bound} and~\ref{thm:tradeoff_bound} together constitute the \emph{quantitative} form of the orthogonality. The propositions below restate the qualitative existence claims that follow as immediate corollaries.

\subsection{Orthogonality Propositions}
\label{sec:orthogonality_props}

\begin{proposition}[RMSE-optimal imputation has positive $F_r$]
\label{prop:rmse_implies_fr}
Let $X_C$ be the block of rows with at least one missing value, and let
$\hat{X}_C^{*} = \arg\min_{\hat{X}}\,\mathbb{E}[\mathrm{RMSE}(\hat{X}, X_C)]$
be the RMSE-minimising imputation.
If the conditional variance $\mathrm{Var}(X_j \mid X_{-j}) > 0$ for any feature $j$ with missing values, then $F_r(\hat{X}_C^{*}) > 0$.
\end{proposition}

\begin{proof}
The RMSE-minimising imputation sets each missing value to its conditional expectation:
\[
  \hat{X}_{ij}^{*} = \mathbb{E}[X_{ij} \mid X_{i,\mathrm{obs}}].
\]
By the law of total variance, for any feature $j$:
\[
  \mathrm{Var}(X_j)
  = \underbrace{\mathbb{E}[\mathrm{Var}(X_j \mid X_{-j})]}_{\text{within-group variance}}
  + \underbrace{\mathrm{Var}(\mathbb{E}[X_j \mid X_{-j}])}_{\text{between-group variance}}.
\]
The conditional-expectation imputation retains only the between-group term.
Since the within-group variance is strictly positive by assumption, the imputed column has strictly lower marginal variance than the true column:
\[
  \mathrm{Var}(\hat{X}_j^{*}) = \mathrm{Var}(\mathbb{E}[X_j \mid X_{-j}])
  < \mathrm{Var}(X_j) \approx \mathrm{Var}(X_j^A),
\]
where $X_j^A$ denotes column $j$ of the available rows $X_A$.
The relative covariance matrix $\tilde{S} = S_A^{-1/2} S_C S_A^{-1/2}$ therefore deviates from the identity, so $D_r(\tilde{S}, I) > 0$, and since $F_r \geq D_r(\tilde{S}, I)$ it follows that $F_r(\hat{X}_C^{*}) > 0$.
\end{proof}

\begin{proposition}[$F_r$-zero imputation has positive RMSE]
\label{prop:fr_implies_rmse}
Let $\hat{X}_C^{**}$ be any imputation satisfying $F_r(\hat{X}_C^{**}) = 0$.
Then $\mathbb{E}[\mathrm{RMSE}(\hat{X}_C^{**}, X_C)] > 0$ unless the true missing values are drawn from the same distribution as $X_A$.
\end{proposition}

\begin{proof}
$F_r = 0$ requires $\hat{X}_C^{**}$ to match $X_A$ in mean, covariance, and skewness --- a constraint on the \emph{marginal distribution} of the imputed rows.
RMSE, by contrast, depends on the \emph{joint distribution} of $(\hat{X}_{ij}^{**},\, X_{ij})$: specifically, $\mathrm{RMSE} = 0$ requires $\hat{X}_{ij}^{**} = X_{ij}$ for every missing cell.
A method can achieve $F_r = 0$ by sampling from $X_A$'s empirical distribution independently of the true missing values, producing $\mathrm{RMSE} > 0$ in expectation.
Equality $\mathrm{RMSE} = 0$ additionally requires each imputed value to equal the corresponding true value, a strictly stronger condition than marginal distributional matching.
Therefore $F_r = 0$ and $\mathrm{RMSE} = 0$ cannot simultaneously hold unless the true $X_C$ is drawn from exactly the same distribution as $X_A$.
\end{proof}

\begin{corollary}[Objectives are non-collinear]
\label{cor:orthogonal}
No imputation method can simultaneously minimise $\mathbb{E}[\mathrm{RMSE}]$ and $F_r$ in general.
Methods that minimise RMSE (e.g.\ MICE) will generically have $F_r > 0$; methods that minimise $F_r$ (e.g.\ MIGA) will generically have RMSE above the irreducible minimum.
\end{corollary}

\begin{remark}
The orthogonality does not imply that $F_r$ and RMSE are \emph{uncorrelated} across methods.
In practice, both tend to decrease as an imputation method improves.
The claim is that \emph{optimising one does not optimise the other}: there is no single imputation that jointly achieves the minimum of each objective.
This is the formal basis for the empirical scatter observed in Chapter~\ref{Chapter4}, Figure~\ref{fig:fr_rmse}.
\end{remark}

The practical consequence is that the choice between MICE and MIGA is not a question of which method is universally better, but which \emph{objective} is appropriate for the downstream application.
When the analysis requires the correct marginal distribution --- density estimation, distributional hypothesis tests, or variance-sensitive statistics --- $F_r$ is the appropriate criterion and MIGA is preferred.
When the goal is accurate reconstruction of individual missing values --- prediction tasks where imputed cells are directly evaluated --- RMSE is the appropriate criterion and MICE or KNN may be preferred.
Chapter~\ref{Chapter4} provides empirical evidence that MIGA's advantage over MICE on $F_r$ is consistent and statistically significant under MAR (Section~\ref{sec:significance}), while MICE retains a RMSE advantage, confirming Corollary~\ref{cor:orthogonal} in practice.

\subsection{Extension to Neural Imputation Methods}
\label{sec:gain_theory}

Propositions~\ref{prop:rmse_implies_fr} and~\ref{prop:fr_implies_rmse} formally apply to any method that either (a) minimises RMSE or (b) achieves $F_r = 0$. This subsection examines how GAIN \cite{Yoon2018} --- a GAN-based neural imputer --- relates to this theory.

\textbf{GAIN's objective.} GAIN trains a generator $G$ to produce imputed values that fool a discriminator $D$, with an additional MSE reconstruction term on observed positions:
\begin{equation}
    \mathcal{L}_G = \underbrace{\mathbb{E}\left[(1-M)\cdot\log D(\hat{X}, H)\right]}_{\text{adversarial}} + \alpha \underbrace{\mathbb{E}\left[M \cdot \|X - G(\tilde{X}, M)\|^2\right]}_{\text{MSE on observed}}
\end{equation}
where $\alpha > 0$ is the reconstruction weight, $M$ is the observed mask, and $H$ is a hint matrix.

\textbf{Interpolation interpretation.} GAIN with $\alpha \to \infty$ forces $G$ to replicate observed values exactly; the discriminator provides no additional signal. In this regime, GAIN effectively reduces to a conditional mean imputer (MICE-like): low RMSE but $F_r > 0$ by Proposition~\ref{prop:rmse_implies_fr}. GAIN with $\alpha \to 0$ removes the MSE constraint; the generator is guided purely by the adversarial term, which pushes toward correct distributional matching. In this regime, GAIN approaches a proper posterior sampler: lower $F_r$ but higher RMSE.

Any finite $\alpha > 0$ therefore places GAIN on the Fr--RMSE Pareto frontier between these extremes. The key claim is:

\begin{corollary}[GAIN cannot achieve $F_r = 0$ and $\mathrm{RMSE} = 0$ simultaneously]
\label{cor:gain_orthogonal}
For any $\alpha > 0$, GAIN's MSE component biases the generator toward conditional means, so $F_r > 0$ by the same law-of-total-variance argument as Proposition~\ref{prop:rmse_implies_fr}. Simultaneously, the adversarial component prevents $G$ from being a pure conditional-mean estimator, so RMSE $> 0$. GAIN therefore cannot jointly achieve $F_r = 0$ and $\mathrm{RMSE} = 0$.
\end{corollary}

\begin{remark}
Corollary~\ref{cor:gain_orthogonal} does \emph{not} claim that GAIN lies on the Fr--RMSE Pareto frontier between MIGA and MICE. At $\alpha = 100$ (the paper's default), experiments (Section~\ref{sec:gain_results}) show GAIN is \emph{dominated} by MICE on Iris: higher $F_r$ (1.478 vs 1.271) \emph{and} higher RMSE (0.124 vs 0.101). The adversarial component adds variance to the imputed values without directing them toward the correct distribution, resulting in off-frontier performance. The corollary's claim is solely the impossibility: no $\alpha > 0$ allows GAIN to achieve both objectives simultaneously.
\end{remark}

This extends the impossibility from the MIGA/MICE pair to any method mixing reconstruction and adversarial objectives. Section~\ref{sec:gain_results} provides empirical confirmation.

%----------------------------------------------------------------------------------------

\section{Genetic Algorithms}

\subsection{Foundations}

A Genetic Algorithm (GA) is a population-based metaheuristic search method inspired by Darwinian natural selection \cite{Holland1975}. A GA maintains a population of candidate solutions, evaluates each candidate using a fitness function, and iteratively produces better solutions by applying selection, crossover (recombination), and mutation operators.

The canonical GA operates as follows \cite{Goldberg1989}:

\begin{enumerate}
    \item \textbf{Initialisation:} Generate an initial population of $l$ individuals, typically at random.
    \item \textbf{Evaluation:} Compute the fitness $f(x)$ for each individual $x$.
    \item \textbf{Selection:} Select individuals for reproduction with probability proportional to fitness.
    \item \textbf{Crossover:} With probability $p_c$, combine pairs of selected individuals to produce offspring.
    \item \textbf{Mutation:} With probability $p_m$, randomly alter components of offspring.
    \item \textbf{Replacement:} Form the new population from survivors and offspring.
    \item \textbf{Repeat} steps 2--6 for $G$ generations or until convergence.
\end{enumerate}

GAs are particularly suited to problems where the fitness landscape is discontinuous, multimodal, or not differentiable --- properties that make gradient-based methods inapplicable.

\subsection{MIGA's Genetic Algorithm}

MIGA \cite{FigueroaGarcia2023} uses a custom GA that departs from the canonical scheme. Most notably, MIGA uses an aggressive diversity injection: at each generation, the majority of the population (approximately 90\% under the paper's default parameters) is replaced with randomly generated individuals drawn from per-variable generators $R_j$. This effectively restarts the search from random each generation, preserving only the best $c$ individuals.

The GA optimises the fitness function $F_r$ over the space of possible imputed values for the missing cells. Each individual in the population represents one complete candidate imputation: a vector of values for all missing cells in the dataset. The fitness function $F_r$ measures distributional distance between the available rows $X_A$ and the completed rows $X_C$, penalising divergence in means, covariance structure, and skewness.

\subsection{Exploration-Exploitation in MIGA}

The paper parameterises the operator outputs as:
\begin{itemize}
    \item $c$: elite individuals carried forward unchanged
    \item $c_1 \times c_3$: mutated copies of elite individuals
    \item $c_2$: crossover offspring from pairs of elites
    \item $l - c - c_1 c_3 - 2(c_2 - 1)$: random individuals from generators (diversity injection)
\end{itemize}

Under the paper's default settings ($l=200$, $c=3$, $c_1=3$, $c_2=2$, $c_3=5$), the diversity injection fills 180/200 = 90\% of the population. This unusual balance means that MIGA is closer to a repeated random restart with elite preservation than a conventional GA.

%----------------------------------------------------------------------------------------

\section{The MIGA Paper}

Figueroa-Garc\'ia, Neruda, and Hernandez-P\'erez (2023) proposed MIGA in ``A genetic algorithm for multivariate missing data imputation,'' published in \textit{Information Sciences} (619, 947--967) \cite{FigueroaGarcia2023}.

\subsection{Problem Formulation}

Given a dataset $X$ with $n$ rows and $p$ features, MIGA partitions it into $X_A$ (the $n_A$ rows with no missing values) and $X_C$ (the $n_C = n - n_A$ rows with at least one missing value). The goal is to find imputed values for the missing cells of $X_C$ such that the statistical properties of the completed $X_C$ match those of $X_A$.

\subsection{The Fitness Function $F_r$}

The fitness function $F_r$ quantifies the distributional distance between $X_A$ and $X_C$ using three terms:

\begin{equation}
    F_r = D_r(\tilde{x}_A,\, \tilde{x}_C) + D_r(\tilde{S},\, I) + D_r(b_A,\, b_C)
    \label{eq:fr_full}
\end{equation}

where $D_r$ denotes the Minkowski distance of order $r$ (typically $r = \infty$, the Chebyshev norm), and:
\begin{itemize}
    \item $\tilde{x} = \bar{x} / S_p$ is the standardised mean vector
    \item $\tilde{S} = S_A^{-1/2} S_C S_A^{-1/2}$ is the relative covariance matrix (equal to the identity when $S_C = S_A$)
    \item $b$ is the bias-corrected skewness vector
\end{itemize}

$F_r = 0$ if and only if $X_C$ has the same mean, covariance, and skewness as $X_A$.

\subsection{Per-Variable Generators}

Each variable $j$ has a random generator $R_j$ used to sample candidate values. The paper specifies that generators are ``obtained from samples'' of the observed values of variable $j$, which this thesis interprets as bootstrap resampling from the empirical distribution (\S\ref{sec:generators}). For integer-valued or discrete variables, this ensures that imputed values are always valid integers.

\subsection{Gaps in the Original Paper}

Three gaps motivate this thesis:

\begin{enumerate}
    \item \textbf{No public implementation.} The paper contains no source code, and implementation details are underspecified at several points.
    \item \textbf{Weak baseline comparison.} CMIM (2005) and ANNI are significantly weaker than modern methods. No comparison against KNN or MICE is provided.
    \item \textbf{No MNAR evaluation.} All experiments assume MAR. The behaviour of MIGA's fitness function under non-ignorable missingness is unknown.
\end{enumerate}

%----------------------------------------------------------------------------------------

\section{Related Work}

\subsection{Distributional Imputation}

The observation that RMSE-optimal imputation suppresses variance has motivated several lines of work. Siddique and Belin (2008) proposed predictive mean matching as a way to preserve distributional properties by imputing observed values rather than model predictions \cite{SiddiqueeBelin2008}. Van Buuren (2018, \S3.4) discusses ``proper imputation'' --- drawing from the full posterior predictive distribution rather than its mean --- as a way to avoid variance suppression \cite{vanBuuren2018}. MIGA can be understood as a non-parametric alternative to proper imputation.

Muzellec et al. (2020) proposed optimal transport-based imputation, which finds imputed values that minimise the Sinkhorn divergence between the completed and observed distributions \cite{Muzellec2020} --- directly analogous to MIGA's $F_r$ objective, but using a differentiable surrogate. They note that distributional methods will generically achieve higher RMSE than regression-based methods, consistent with the findings of this thesis.

\subsection{Neural Network Imputation}

GAIN \cite{Yoon2018} uses a Generative Adversarial Network to generate imputations, with a discriminator that learns to distinguish imputed from observed values --- conceptually similar to MIGA's distributional objective but in a learned feature space. GRAPE \cite{You2020} formulates imputation as a graph completion problem using message passing over a bipartite graph of observations and features. A comparison against GAIN or GRAPE is identified as future work (\S\ref{sec:future_neural}).

\subsection{Evaluation of Imputation Methods}

Oberman and Vink (2024) survey evaluation practices for imputation and find no universal standard for RMSE normalisation, making cross-paper comparisons unreliable \cite{ObermanVink2024}. This thesis uses range-normalised RMSE to match the paper's reported baseline values. Ipsen et al. (2022) show that evaluating imputation quality under MNAR using RMSE on the missing cells is itself biased, because the missing cells are a non-representative sample of the marginal distribution \cite{Ipsen2022} --- a point directly relevant to Chapter~\ref{Chapter4}'s MNAR analysis.
